%% マニュアル作成要件の追補（2026-10-24）で確かめた試作。書式の例（bib の枠＋出力）と段落の体裁。
%% 組版：TEXINPUTS=".:<repo>//:" lualatex → biber → lualatex ×2（examples/ と 2026-10-24-manual-proto-ex.lua を同じフォルダに置く）
\documentclass[a4paper,10pt,oneside,line_length=37zw,number_of_lines=42]{jlreq}
\usepackage[american,japanese]{babel}
\usepackage[autostyle=true]{csquotes}
\usepackage{xurl}
\usepackage[haranoaji]{luatexja-preset}
\usepackage{fontspec}
\setmainfont{Libertinus Serif}\setsansfont{Libertinus Sans}\setmonofont{Latin Modern Mono}[Scale=0.9]
\usepackage{btxdockit}
\renewcommand*{\marglistfont}{\spotcolor\sffamily\small}
\renewcommand*{\optionlistfont}{\spotcolor\sffamily\displayverbfont}
\renewcommand*{\ltxsyntaxlabelfont}{\spotcolor\displayverbfont}
\hypersetup{colorlinks,linkcolor=spot,urlcolor=spot}
\ModifyHeading{section}{font=\sffamily\gtfamily\bfseries\spotcolor\Large}
\ModifyHeading{subsection}{font=\sffamily\gtfamily\bfseries\spotcolor\large}
\usepackage[backend=biber,style=jsr,date=ymd]{biblatex}
\addbibresource{examples/names.bib}\addbibresource{examples/trans.bib}
%% 段落：字下げなし，段落間を少し空ける
\setlength{\parindent}{0pt}
\setlength{\parskip}{.5\baselineskip plus .1\baselineskip}
%% bib の枠（biblatex の style=bibtex：白地・細い黒枠）
\lstdefinestyle{jsrbib}{basicstyle=\ttfamily\small\setlength{\parskip}{0pt},frame=single,framerule=0.4pt,rulecolor=\color{black},
  framesep=3pt,xleftmargin=3.4pt,xrightmargin=3.4pt,aboveskip=.6\baselineskip,belowskip=.3\baselineskip,
  columns=flexible,keepspaces=true,breaklines=true,showstringspaces=false,
  moredelim=[s][\itshape\rmfamily]{<}{>},
  moredelim=[l][\color{gray}\rmfamily\footnotesize]{※}}
\makeatletter\lst@AddToHook{Init}{\parskip\z@skip}\lst@AddToHook{EveryLine}{\parskip\z@skip}\lst@AddToHook{EveryPar}{\parskip\z@skip}\makeatother
\lstnewenvironment{bibtemplate}{\setlength{\parskip}{0pt}\lstset{style=jsrbib}}{}
\directlua{dofile('2026-10-24-manual-proto-ex.lua')}
\makeatletter
\newcounter{jsrex}
%% 例ごとに refsection を分け，指定したキーだけを \nocite して文献リストを組む
\newcommand*{\bibexample}[1]{%
  \stepcounter{jsrex}%
  \directlua{jsrex.extract("names,trans", "\luaescapestring{#1}", "\jobname-ex\thejsrex.tmp")}%
  \begingroup\setlength{\parskip}{0pt}\lstinputlisting[style=jsrbib]{\jobname-ex\thejsrex.tmp}\endgroup
  \begin{refsection}\nocite{#1}%
  \begingroup\small\setlength{\parskip}{0pt}%
  \setlength{\marginparwidth}{\marglistwidth}\setlength{\marginparsep}{\marglistsep}\reversemarginpar
  \def\jsr@first{1}%
  \renewbibmacro*{begentry}{\ifdefstring\jsr@first{1}{\marginpar{\raggedleft\marglistfont 出力}\gdef\jsr@first{0}}{}}%
  \printbibliography[heading=none]%
  \endgroup
  \end{refsection}\par}
\makeatother
\begin{document}
\section{書式と例}
\subsection{単著・共著（\bibtype{book}）}
著者が1人の書籍は，\bibtype{book} に書きます。和文の著者には，並べ替え用の \bibfield{sortname}（ローマ字）を添えます。

記入のひな形は次のとおりです：
\begin{bibtemplate}
@book{<キー>,
  author    = {<姓>,<名>},       ※ 和文は姓と名を「,」で区切る
  sortname  = {<Family>, <Given>}, ※ 和文のみ。ローマ字
  date      = {<刊行年>},
  title     = {<書名>},
  publisher = {<出版社>},
  langid    = {japanese},
}
\end{bibtemplate}
たとえば，次のように書きます：
\bibexample{aoki2015}
共著は，著者を \texttt{and} でつなぎます。文献リストでは中黒でつながります：
\bibexample{aokikita2018}
欧文の共著は，2人目以降を名・姓の順に出し，最後を and でつなぎます：
\bibexample{clark2017}
\subsection{翻訳書（\bibfield{related}）}
訳書は原典のエントリを作り，\bibfield{related} でつなぎます。文献リストには，原典の項目に訳書が添えられて1項目で出ます：
\bibexample{hale1962,hale1985}
原典を引用していなくても，原典のエントリを bib に書いておけば，同じ形で出ます（\bibfield{related} が指す原典が代表になります）：
\bibexample{hale1985}
\section{コマンドの例}
\begin{ltxsyntax}
\cmditem{parencite}[prenote][postnote]{key}
\cmditem{textcite}[prenote][postnote]{key}

全角の丸括弧で囲む文献注です。本文での引用の基本形です。ページは \prm{postnote} に書きます。
\end{ltxsyntax}
\begin{fieldlist}
\fielditem{origpublisher}{literal}

紙媒体の出版社。電子書籍で使います。

二つ目の段落です。段落頭は下げず，段落の間を少し空けます。
\end{fieldlist}
\end{document}
